\scanpage{15}
canoniquement aux surfaces riemanniennes standart \uline{types} ${\unsure{\mathbb S|S}}^2,\mathbb R^2=\mathbb C,\mathbb H$ avec le \textquotedblleft répère\textquotedblright\ différentiel habituel \ldots. On utilise \unsure{comme|encore} notation $C_{p,q}$

\deleted{\unsure{6}} On distingue \deleted{types} des cartes standard de type sphérique, euclidien ou hyperbolique.

\textcircled{6} \textcircled{a} \deleted{\unsure{p,q}} $C_{p,q}$ est sphérique \deleted{\unsure{sss}} exactement dans les cas suivants : \deleted{\unsure{$p=2,q$|$p=q$}} $p=2$, (\deleted{\unsure{$q$ quelconque $\geq1$}}\inserted{$1\leq q<+\infty$}) ; $q=2$ (\deleted{\unsure{$p$ quelconque $\geq1$}}\inserted{$1\leq p<+\infty$}), \deleted{$p=q=3,\ p=$}
\diagram{\(\begin{array}{ccccc}
(3,3)&(3,4)&(4,3)&(3,5)&(5,3)\\
\text{tétraèdre}&\text{cube}&\text{octaèdre}&\text{dodécaèdre}&\text{icosaèdre}
\end{array}\)}
-- ou encore sss $C_{p,q}$ \uline{finie} ou encore sss $\Gamma_{p,q}$ fini.

\textcircled{b} $C_{p,q}$ est euclidien exactement dans les cas suivants : \inserted{$(2,\infty)$, \deleted{\unsure{$(\infty,$|$(\infty,2)$}} $(\infty,2)$}
\diagram{\(\begin{array}{ccc}
(4,4),&(3,6),&(6,3)\\
\text{pavage carré}&\text{pavage}&\text{pavage}\\
&\text{hexagonal}&\text{triangulaire}
\end{array}\)}
\deleted{\unsure{$(2,\infty)$}} -- ou encore sss $\Gamma_{pq}$ est infini et \uline{résoluble}.

\textcircled{c} $C_{pq}$ est hyperbolique dans les autres cas i.e. \deleted{\unsure{$(p,q)>(3,3)$}} \deleted{\unsure{$(p>3,\ q\geq3)$|$(p\geq3,\ q\geq3)$}} \deleted{\unsure{ou $(p=3$,}}
\[
\text{\textcircled{$c_1$}}\quad
\left\{\begin{aligned}
&\text{ou }\deleted{\unsure{(p\geq3,q\geq3,\ (p,q)\ne(4,4))\ \text{ou}}}\qquad(p,q)>(4,4)\\
&\text{ou }\quad p=3,\ q\geq7\\
&\text{ou }\quad q=3,\ p\geq7
\end{aligned}\right.
\]
\margin{ordre produit sur $\unsure{\overline{\mathbb N}^{\,*}\times\overline{\mathbb N}^{\,*}|\mathbb N^*\times\overline{\mathbb N}^{\,*}}$}
\[
\begin{aligned}
\text{\textcircled{$c'_1$}}\ &\text{ou encore }\quad p+q\geq10\ \text{ou }(p,q)=(4,5)\ \text{ou }(p,q)=(5,4)\\
\unsure{(c''_1)|(c'_1)}\ &\text{ou encore }\quad p+q\geq9\ \text{et }(p,q)\ne(3,6),(6,3)\quad\deleted{\unsure{\text{grpe}}}\\
(c_2)\ &\text{ou encore }\quad\Gamma_{pq}\ \text{infini non résoluble}\quad\deleted{\unsure{(?)}}
\end{aligned}
\]
$(c'_2)$ ou encore $\Gamma_{pq}$ contient ss-groupe libre à 2 générateurs (donc contient \unsure{un} groupe libre à $n$ générateurs $\forall n\geq0$)

\margin{$\checkmark$}
\textcircled{7} Les cartes \inserted{quasi-}régulières \inserted{$(X\supset K\supset S)$} épinglées \inserted{de type $(p,q)$} s'obtiennent (à isom canonique près dans la cat. isotopique des cartes) à partir des sous-groupes $\pi\subset\Gamma_{p,q}$
% archive sheet number: 14

